\subsection{\texorpdfstring{SYSTEMS OF TYPE \(D_{l}\) (\(l \geqslant 3\))}{SYSTEMS OF TYPE D\_l (l >= 3)}}

\begin{enumerate}
\def\labelenumi{(\Roman{enumi})}
\item
  Consider in \(V = R^{l}\) the group \(L_{0}\) (no. 4). The set R of \(\alpha \in L_{0}\) such that \((\alpha|\alpha) = 2\) consists of the vectors \(\pm\varepsilon_{i} \pm \varepsilon_{j}\) ( \(1 \leqslant i < j \leqslant l\) ). It is clear that R generates V and that \(2(\alpha|\beta)/(\alpha|\alpha) \in \mathbf{Z}\) for all \(\alpha, \beta \in R\) . Thus R is a reduced root system in V (no. 4). The number of roots is \(n = 2l(l - 1)\) .
\item
  Put
\end{enumerate}

\[
\alpha_ {1} = \varepsilon_ {1} - \varepsilon_ {2}, \alpha_ {2} = \varepsilon_ {2} - \varepsilon_ {3} \dots . \alpha_ {l - 1} = \varepsilon_ {l - 1} - \varepsilon_ {l}, \alpha_ {l} = \varepsilon_ {l - 1} + \varepsilon_ {l}.
\]

The following formulas are immediate:

\[
\begin{array}{c} \varepsilon_ {i} - \varepsilon_ {j} = \alpha_ {i} + \alpha_ {i + 1} + \dots + \alpha_ {j - 1} (i <   j) \\ \varepsilon_ {i} + \varepsilon_ {j} = \alpha_ {i} + \alpha_ {i + 1} + \dots + \alpha_ {j - 1} + 2 \alpha_ {j} + 2 \alpha_ {j + 1} + \dots \\ \qquad \qquad \qquad \qquad \qquad \qquad \dots + 2 \alpha_ {l - 2} + \alpha_ {l - 1} + \alpha_ {l} (i <   j \leqslant l - 2) \\ \varepsilon_ {i} + \varepsilon_ {l - 1} = \alpha_ {i} + \alpha_ {i + 1} + \dots + \alpha_ {l} (i <   l - 1) \\ \varepsilon_ {i} + \varepsilon_ {l} = \alpha_ {i} + \alpha_ {i + 1} + \dots + \alpha_ {l - 2} + \alpha_ {l} (i <   l - 1) \\ \varepsilon_ {l - 1} + \varepsilon_ {l} = \alpha_ {l}. \end{array}
\]

so \((\alpha_{1},\ldots,\alpha_{l})\) is a basis if R (§ 1, no. 2. Cor. 3 of Prop. 20). Further, \(\| \alpha_{i}\|^{2}=2\) for all \(i\) , \((\alpha_{i}|\alpha_{j})=0\) for \(i+1<j\) except for \(i=l-2\) , \(j=l\) in which case \((\alpha_{l-2}|\alpha_{l})=-1\) , \((\alpha_{i}|\alpha_{i+1})=-1\) for \(i\leqslant l-2\) , and finally \((\alpha_{l-1}|\alpha_{l})=-1\) ; the Dynkin graph of R is thus of type D \(_{l}\) . The positive roots are the \(\varepsilon_{i}\pm\varepsilon_{j}\) for \(i<j\) .

\begin{enumerate}
\def\labelenumi{(\Roman{enumi})}
\setcounter{enumi}{2}
\item
  We have \(h = n / l = 2(l - 1)\) .
\item
  Let \(\tilde{\alpha} = \varepsilon_{1} + \varepsilon_{2} = \alpha_{1} + 2\alpha_{2} + \cdots + 2\alpha_{l-2} + \alpha_{l-1} + \alpha_{l}\) , which is a root. The sum of its coordinates relative to \((\alpha_{i})\) is
\end{enumerate}

\[
2 l - 3 = h - 1.
\]

Hence \(\tilde{\alpha}\) is the highest root.

If \(l = 3\) , we have

\[
(\tilde {\alpha} | \alpha_ {1}) = 0, \quad (\tilde {\alpha} | \alpha_ {2}) = (\tilde {\alpha} | \alpha_ {3}) = 1.
\]

If \(l \geqslant 4\) . we have \((\tilde{\alpha}|\alpha_i) = 0\) for \(i \neq 2\) and \((\tilde{\alpha}|\alpha_2) = 1\) . Hence the completed Dynkin graph is:

\[
\alpha_ {1} \quad \alpha_ {2} \quad \alpha_ {3} \quad \dots \quad \alpha_ {i - 3} \quad \alpha_ {i - 2} \quad \alpha_ {i - 1}
\]

\begin{enumerate}
\def\labelenumi{(\Alph{enumi})}
\setcounter{enumi}{21}
\tightlist
\item
  Since \((\alpha |\alpha) = 2\) for all \(\alpha \notin R\) , we have \(R^{\sim} = R\) .
\end{enumerate}

The length of the roots for \(\Phi_{\mathrm{R}}\) is \(h^{-1 / 2} = (2l - 2)^{-1 / 2}\) . Hence

\[
\Phi_ {R} (x, y) = (x | y) / (4 l - 4) \quad \text { and } \quad \gamma (R) = 4 (l - 1) ^ {2}.
\]

\begin{enumerate}
\def\labelenumi{(\Roman{enumi})}
\setcounter{enumi}{5}
\tightlist
\item
  A calculation analogous to that in no. 7 gives the fundamental weights:
\end{enumerate}

\[
\begin{array}{r l} \omega_ {i} & = \varepsilon_ {1} + \varepsilon_ {2} + \dots + \varepsilon_ {i} \\ & = \alpha_ {1} + 2 \alpha_ {2} + \dots + (i - 1) \alpha_ {i - 1} + i (\alpha_ {i} + \alpha_ {i + 1} + \dots + \alpha_ {l - 2}) \\ & \quad + \frac {1}{2} i (\alpha_ {l - 1} + \alpha_ {l}) \end{array}
\]

for i \textless{} l - 1,

\[
\begin{array}{l} \omega_ {l - 1} = \frac {1}{2} (\varepsilon_ {1} + \varepsilon_ {2} + \dots + \varepsilon_ {l - 2} + \varepsilon_ {l - 1} - \varepsilon_ {l}) \\ \qquad = \frac {1}{2} (\alpha_ {1} + 2 \alpha_ {2} + \dots + (l - 2) \alpha_ {l - 2} + \frac {1}{2} l \alpha_ {l - 1} + \frac {1}{2} (l - 2) \alpha_ {l}), \\ \omega_ {l} = \frac {1}{2} (\varepsilon_ {1} + \varepsilon_ {2} + \dots + \varepsilon_ {l - 2} + \varepsilon_ {l - 1} + \varepsilon_ {l}) \\ \qquad = \frac {1}{2} (\alpha_ {1} + 2 \alpha_ {2} + \dots + (l - 2) \alpha_ {l - 2} + \frac {1}{2} (l - 2) \alpha_ {l - 1} + \frac {1}{2} l \alpha_ {l}). \end{array}
\]

\begin{enumerate}
\def\labelenumi{(\Roman{enumi})}
\setcounter{enumi}{6}
\tightlist
\item
  The sum of the positive roots is
\end{enumerate}

\[
\begin{array}{l} 2 \rho = 2 (l - 1) \varepsilon_ {1} + 2 (l - 2) \varepsilon_ {2} + \dots + 2 \varepsilon_ {l - 1} \\ = \sum_ {i = 1} ^ {l - 2} 2 (i l - \frac {i (i + 1)}{2}) \alpha_ {i} + \frac {l (l - 1)}{2} (\alpha_ {l - 1} + \alpha_ {l}). \end{array}
\]

\begin{enumerate}
\def\labelenumi{(\Roman{enumi})}
\setcounter{enumi}{7}
\item
  The \(\pm \varepsilon_{i} \pm \varepsilon_{j}\) generate \(L_{1}\) (no. 4), so \(Q(R) = L_{1}\) . Hence \(Q(R^{-}) = L_{1}\) and consequently \(P(R) = L_{2}\) (no. 4). By no. 4, \(P(R)/Q(R)\) is isomorphic to \(\mathbf{Z}/4\mathbf{Z}\) for \(l\) odd, and to \(\mathbf{Z}/2\mathbf{Z} \times \mathbf{Z}/2\mathbf{Z}\) for \(l\) even. In the first case, \(P(R)/Q(R)\) is generated by the canonical image of \(\omega_{l}\) (and also by that of \(\omega_{l-1}\) ). In the second case, \(P(R)/Q(R)\) is generated by the canonical images of \(\omega_{l-1}\) and \(\omega_{l}\) . In both cases the connection index is 4.
\item
  and (X) In \(\mathbf{R}^l\) , the orthogonal reflection \(s_{\varepsilon_i - \varepsilon_j}\) ( \(i \neq j\) ) interchanges \(\varepsilon_i\) and \(\varepsilon_j\) and leaves invariant the \(\varepsilon_k\) with \(k\) distinct from \(i\) and \(j\) . The \(s_{\varepsilon_i - \varepsilon_j}\) generate a group \(\mathbf{G}_1\) isomorphic to the symmetric group \(\mathfrak{S}_l\) . On the other hand, \(s_{ij} = s_{\varepsilon_i - \varepsilon_j} s_{\varepsilon_i + \varepsilon_j}\) transforms \(\varepsilon_i\) to \(-\varepsilon_i\) , \(\varepsilon_j\) to \(-\varepsilon_j\) and leaves invariant the \(\varepsilon_k\) with \(k\) distinct from \(i\) and \(j\) . The \(s_{ij}\) generate a group \(\mathbf{G}_2\) , the set of automorphisms \(u\) of the vector space \(\mathbf{R}^l\) such that \(u(\varepsilon_i) = (-1)^{\nu_i} \varepsilon_i\) with \(\prod_{i=1}^{l} (-1)^{\nu_i} = 1\) . The group \(\mathbf{G}_2\) is isomorphic to \((\mathbf{Z}/2\mathbf{Z})^{l-1}\) , and \(\mathbf{G}_2\) is normal in W(R), so W(R) is isomorphic to the semi-direct product of \(\mathfrak{S}_l\) by \((\mathbf{Z}/2\mathbf{Z})^{l-1}\) . Consequently, its order is \(2^{l-1}.l!\) .
\end{enumerate}

The polynomial functions \(t_j\) of no. 5 are invariant under \(W(R)\) , and so is \(t = \xi_1\xi_2\ldots\xi_l\) ; moreover, \(t_l = t^2\) . Let \(P(\xi_1,\ldots ,\xi_l)\) be a polynomial invariant under \(W(R)\) . Let \(\xi_1^{\nu_1}\xi_2^{\nu_2}\ldots\xi_l^{\nu_l}\) be a monomial featuring in \(P\) such that \(\nu_{i}\) is odd; then \(\nu_{j}\) is odd for all \(j\) because the monomial \((-1)^{\nu_i + \nu_j}\xi_1^{\nu_1}\xi_2^{\nu_2}\ldots\xi_l^{\nu_l}\) features in \(s_{ij}(P)\) , so \(\nu_{i} + \nu_{j} \equiv 0 \pmod{2}\) and \(\nu_{j} \equiv 1 \pmod{2}\) . Thus

\(P = P_{1} + tP_{2}\) , where all the monomials featuring in \(P_{1}\) and \(P_{2}\) have only even exponents. Since P is invariant under the permutations of the \(\xi_{i}\) . \(P_{1}\) and \(P_{2}\) have the same property, and so can be written as polynomials in \(t_{1}, t_{2}, \ldots, t_{l}\) . This proves that the algebra \(\mathbf{S}(\mathbf{R}^{l})^{\mathrm{W}(\mathbb{R})}\) is generated by \(t_{1}, t_{2}, \ldots, t_{l-1}, t\) . Moreover, the transcendence degree of the field of fractions of \(\mathbf{S}(\mathbf{R}^{l})^{\mathrm{W}(\mathbb{R})}\) is l, so \(t_{1}, t_{2}, \ldots, t_{l-1}, t\) are algebraically independent. We conclude that the sequence of exponents, suitably ordered, is:

\[
1, 3, 5, \dots , 2 l - 5, 2 l - 3, l - 1.
\]

Note that \(l - 1\) appears twice if \(l\) is even, and once if \(l\) is odd.

\begin{enumerate}
\def\labelenumi{(\Roman{enumi})}
\setcounter{enumi}{10}
\tightlist
\item
  The automorphisms of the Dynkin graph are those of the underlying graph. Thus:
\end{enumerate}

\begin{enumerate}
\def\labelenumi{\arabic{enumi})}
\item
  If \(l = 3\) . A(R)/W(R) is isomorphic to Z/2Z.
\item
  If \(l = 4\) , every permutation of the terminal vertices defines an automorphism of the graph, so \(\mathrm{A}(\mathrm{R}) / \mathrm{W}(\mathrm{R})\) is isomorphic to \(\mathfrak{S}_3\) .
\item
  If \(l \geqslant 5\) , the chains starting at the ramification point have length 1,1, and \(l - 3 \geqslant 2\) . The only automorphism of the graph distinct from the identity thus corresponds to the automorphism \(\varepsilon \in \mathrm{A}(\mathrm{R})\) which interchanges \(\alpha_{l-1}\) and \(\alpha_l\) and leaves fixed the \(\alpha_i\) for \(1 \leqslant i \leqslant l - 2\) . Thus \(\mathrm{A}(\mathrm{R}) / \mathrm{W}(\mathrm{R})\) is isomorphic to \(\mathbf{Z}/2\mathbf{Z}\) ; moreover, \(\mathrm{A}(\mathrm{R})\) is the semi-direct product of the group \(\mathbf{G}_1 \cong \mathfrak{S}_l\) defined in (IX) by the group \(\mathbf{G}_3\) consisting of the automorphisms \(u\) of \(\mathbf{R}^l\) such that \(u(\varepsilon_i) = \pm \varepsilon_i\) for all \(i\) .
\end{enumerate}

If \(l\) is even, \(-1 \in \mathrm{W}(\mathrm{R})\) , so \(w_0 = -1\) . If \(l\) is odd, \(-1 \notin \mathrm{W}(\mathrm{R})\) , so \(\mathrm{A}(\mathrm{R}) = \mathrm{W}(\mathrm{R}) \times \{1, -1\}\) and \(w_0 = -\varepsilon\) .

\begin{enumerate}
\def\labelenumi{(\Roman{enumi})}
\setcounter{enumi}{11}
\tightlist
\item
  For \(l\) even, \(\mathrm{P}(\mathbb{R}^{-}) / \mathrm{Q}(\mathbb{R}^{-})\) has three elements of order 2, namely \(\omega_{1}, \omega_{l-1}\) and \(\omega_{l}\) . Since \(\omega_{l}\) (resp. \(\omega_{l-1}\) ) interchanges the vertices corresponding to \(\alpha_{0}\) and \(\alpha_{l}\) (resp. \(\alpha_{l-1}\) ), it interchanges those corresponding to \(\alpha_{1}\) and \(\alpha_{l-1}\) (resp. \(\alpha_{l}\) ) and also those corresponding to \(\alpha_{j}\) and \(\alpha_{l-j}\) for
\end{enumerate}

\[
2 \leqslant j \leqslant l - 2.
\]

We have \(\omega_{1} = \omega_{l}\omega_{l - 1}\)

For \(l\) odd, \(\mathrm{P}(\mathbb{R}^{-}) / \mathrm{Q}(\mathbb{R}^{-})\) has two elements of order 4, namely \(\omega_{l-1}\) and \(\omega_l\) , and one element of order 2, equal to \(\omega_1\) . Indeed, \(\omega_1\) interchanges the vertices corresponding to \(\alpha_0\) and \(\alpha_1\) , so it leaves fixed the vertices corresponding to \(\alpha_j\) for \(2 \leqslant j \leqslant l-2\) and is necessarily of order 2. Consequently \(\omega_l\) is of order 4 and transforms the vertex corresponding to \(\alpha_0\) (resp. \(\alpha_l\) , resp. \(\alpha_1\) , resp. \(\alpha_{l-1}\) ) to that corresponding to \(\alpha_l\) (resp. \(\alpha_1\) , resp. \(\alpha_{l-1}\) , resp. \(\alpha_0\) ), and interchanges the vertices corresponding to \(\alpha_j\) and \(\alpha_{l-j}\) for \(2 \leqslant j \leqslant l-2\) . We have \(\omega_1 = \omega_l^2\) and \(\omega_{l-1} = \omega_l^3\) .

For \(l \neq 4\) , the non-identity element of \(\mathrm{A}(\mathrm{R})/\mathrm{W}(\mathrm{R})\) interchanges the vertices corresponding to \(\alpha_{l-1}\) and \(\alpha_{l}\) , and consequently interchanges the elements \(\omega_{l-1}\) and \(\omega_{l}\) of \(\mathrm{P}(\mathrm{R}^{-})/\mathrm{Q}(\mathrm{R}^{-})\) . For l odd, the automorphism of \(\mathrm{P}(\mathrm{R}^{-})/\mathrm{Q}(\mathrm{R}^{-})\) thus obtained is the map \(x \mapsto -x\) .

For \(l = 4\) , \(A(R)/W(R)\) can be identified with the group of permutations of \(\{1,3,4\}\) and acts by permutations of the indices on \(\{\omega_1,\omega_3,\omega_4\}\) .

